% !TEX root = ../../main.tex

\section{系统总体架构设计}

HearBridge 系统采用分层式多端协同架构，整体可划分为用户访问层、业务服务层、智能识别与语义服务层以及数据与资源支撑层。
用户访问层包括 HarmonyOS 移动端和 Vue Web 管理端；
业务服务层以 Spring Boot 服务端为核心；
智能识别与语义服务层包括 Python 手势识别服务和 DeepSeek 语义修正服务；
数据与资源支撑层包括 MySQL、Redis 和 MinIO 等组件。

系统总体架构如图~\ref{fig:chapter04-system-architecture} 所示。

\WhuImageFigure[0.95\textwidth][0.62\textheight]{figures/chapter04/fig4-1-system-architecture.png}{系统总体架构图}{fig:chapter04-system-architecture}

如图~\ref{fig:chapter04-system-architecture} 所示，HarmonyOS 移动端主要面向普通用户，提供手语资源浏览、动作示范展示、实时手势训练、句子视频识别和个人中心等功能。
移动端通过 HTTP 请求访问 Spring Boot 服务端，完成用户认证、资源查询、视频识别请求提交和结果获取等操作；
在实时手势识别场景中，移动端通过 WebSocket 与 Python 手势识别服务进行通信，以支持连续帧传输和实时识别结果反馈。

Vue Web 管理端主要面向系统管理员，提供手势分类管理、手势资源管理、样本管理、模型训练管理、模型版本管理和系统概览等功能。
管理端通过 HTTP 请求访问 Spring Boot 服务端，由服务端完成权限校验、数据读写、文件资源管理和识别服务调用。

Spring Boot 服务端是系统业务服务中心，主要负责用户认证与权限校验、手势分类管理、手势资源管理、样本管理、模型版本管理、文件资源管理、句子视频识别请求编排和 DeepSeek 语义修正调用。
服务端向下访问 MySQL、Redis 和 MinIO 等数据与资源支撑组件，并根据业务需要调用 Python 手势识别服务和 DeepSeek 语义修正服务。

Python 手势识别服务主要负责实时手势识别、句子视频识别、样本处理和模型训练与加载等功能。
该服务不直接承担用户业务管理职责，而是作为独立识别能力服务被移动端或服务端调用。
DeepSeek 语义修正服务主要用于对句子视频识别结果进行语义整理和自然化表达，使最终展示结果更加便于用户理解。

数据与资源支撑层中，MySQL 用于存储用户、管理员、手势分类、手势资源、样本和模型版本等结构化业务数据；
Redis 用于存储登录状态和临时缓存数据；
MinIO 用于存储头像、手势资源封面、动作示范资源、样本视频和模型文件等非结构化资源。
通过这种分层设计，系统能够在保证业务功能完整性的同时，提高后续维护和扩展能力。
